\section{Limitations}
\label{sec:limitations}

This survey has several limitations that qualify our conclusions.

\textbf{Recency bias.} The field evolves rapidly; new systems appear
weekly. Our survey covers publications and releases through February
2026. Systems published after this date are not included.

\textbf{Selection bias.} We identify systems primarily through
English-language academic databases, GitHub, and practitioner blogs.
Systems published in non-English venues or deployed in closed-source
products may be underrepresented.

\textbf{Classification subjectivity.} Mapping systems to the CoALA
taxonomy, RAG categories, and lifecycle operations requires
interpretive judgment. We mitigate this through explicit definitions
(\Cref{sec:taxonomy}) and consistent application, but reasonable
disagreements are possible.

\textbf{Benchmark comparability.} Not all systems report results on
the same benchmarks. LoCoMo, LongMemEval, and DMR scores are not
directly comparable across papers due to variations in evaluation
setup, model choice, and preprocessing.

\textbf{Reproducibility.} All baseline experiments
(\S\ref{sec:baseline-replication}, \S\ref{sec:mab-replication}) use
Gemini-2.0-Flash-001 via OpenRouter (accessed February 2026;
temperature~0, max\_tokens~512) as the generator and GPT-4o-mini via
OpenRouter (temperature~0, max\_tokens~10) as the judge.  Embedding
model: \texttt{all-MiniLM-L6-v2} via ChromaDB.  Total inference cost
for all three benchmarks: approximately \$12.  Code and raw results are
available at~\cite{alaya2026}.  Model
behavior may change with API updates; our results reflect model versions
available in February 2026.

\textbf{No system-level evaluation.} This paper surveys and classifies
memory architectures but does not evaluate any specific system's
end-to-end performance. Comparative evaluation of lifecycle operations
(forgetting, consolidation, transformation) across systems remains
future work.
